% Clinical Background section. Paste into your own thesis, or \input it.
%
% Needs graphicx. Figures:
%   figures/anatomy_ct.png       <- python scripts/make_anatomy_figure.py
%   figures/external/*.pdf|jpg   <- python scripts/fetch_external_figures.py
% Cohort numbers here            <- python scripts/cohort_numbers.py
% Credits for the borrowed figures are generated into figures/external/CREDITS.md.
% The captions below must match that file; both come from the same Commons API response.

\section{Heart anatomy}
\label{sec:anatomy}

The heart drives the circulation of blood that delivers oxygen and nutrients to every tissue in the
body and carries away metabolic waste \cite{shah2009heart}.
It is a muscular organ that contracts rhythmically throughout life and lies in the mediastinum, the
compartment of the chest between the lungs \cite{weinhaus2024anatomy}.
It has four chambers: two thin-walled \emph{atria} that receive blood, the right atrium (RA) and the
left atrium (LA), and two thicker \emph{ventricles} that pump it out, the right ventricle (RV) and
the left ventricle (LV) \cite{weinhaus2024anatomy}.

The right and left sides operate as two pumps in series, and the path blood takes between them is
shown in Figure~\ref{fig:heart-chambers}.
Blood that has given up its oxygen to the tissues returns through the superior and inferior venae
cavae, the body's two largest veins, into the RA, which passes it into the RV
\cite{weinhaus2024anatomy}.
The RV pumps it through the lungs, where it is oxygenated --- this loop is the \emph{pulmonary
circuit} --- and it returns to the LA, which passes it into the LV.
The LV, the thickest-walled chamber, ejects it into the aorta, the body's largest artery, which
begins the \emph{systemic circuit} that supplies the rest of the body \cite{weinhaus2024anatomy}.
Blood returning from those tissues enters the venae cavae again, and the cycle repeats.

\begin{figure}[t]
  \centering
  \includegraphics[width=0.7\textwidth]{external/heart_chambers.pdf}
  \caption{Blood crosses the heart twice, and the coronary arteries branch off at the exit. The
    four chambers, four valves and the direction of flow; the coronary arteries are not drawn, and
    arise from the aorta immediately above the aortic valve at the top of the figure. Figure by
    Wapcaplet, CC BY-SA 3.0, via Wikimedia Commons.}
  \label{fig:heart-chambers}
\end{figure}

The heart's own muscle is not supplied by the blood within the chambers it pumps.
A separate set of arteries, the coronary arteries, leaves the aorta and runs over the outside of the
heart, feeding the muscle from there \cite{devos2024coronary}.
That separate supply is what makes coronary artery disease possible.

\subsection{The origin of the coronary arteries}
\label{sec:chambers}

Because the two circuits run in series through one pump, unidirectional flow is maintained by four
valves, one at the inlet and one at the outlet of each ventricle
\cite{weinhaus2024anatomy,shah2009heart}.
One of them matters here beyond that mechanical role.
The \emph{aortic valve} sits at the exit of the LV, and the coronary arteries arise from the wall of
the aorta immediately above it \cite{devos2024coronary}, so the arterial tree studied in this thesis
begins at the point where the LV ejects.

\subsection{Coronary perfusion}
\label{sec:perfusion}

The heart alternates between two phases, and the coronary arteries are affected differently by each.
In \emph{diastole} a chamber relaxes and fills with blood; in \emph{systole} it contracts and ejects
what it holds \cite{shah2009heart}.
Atria and ventricles take turns, so the ventricles are in diastole while the atria are in systole.

The coronary circulation is unusual in that its supply is interrupted by the organ it supplies.
The heart wall has three layers, shown in Figure~\ref{fig:heart-wall}: a thin inner lining, the
\emph{endocardium}; a thick muscular middle layer, the \emph{myocardium}, which does the mechanical
work of pumping; and a thin outer covering, the \emph{epicardium}, beneath which the coronary
arteries run.
Those arteries send branches down into the myocardium, so contraction compresses them: perfusion,
the delivery of blood to the tissue, is impeded during systole, and blood is briefly forced back
out of the compressed vessels \cite{duncker2015regulation}.
During diastole the muscle relaxes, inflow rises, and it rises most in the \emph{subendocardium},
the inner third of the muscle and the layer that was compressed most \cite{duncker2015regulation}.
Coronary perfusion is therefore predominantly a diastolic event, and a fast heart rate, by
shortening diastole, reduces the time available for it \cite{duncker2015regulation}.

The heart already extracts 60--80\% of the oxygen carried past it at rest, so additional demand
must be met by increased flow rather than by increased extraction \cite{duncker2015regulation}.
A vessel that cannot supply that increase produces \emph{ischemia}, muscle working without enough
oxygen, which is the mechanism by which a narrowing in a coronary artery becomes symptomatic
\cite{duncker2015regulation}.
The margin is narrowest in the subendocardium: it extracts 83\% of the oxygen carried past it
against 70\% in the \emph{subepicardium}, the outer third of the muscle, and it is also the layer
most compressed during systole \cite{duncker2015regulation}.

This decides how the arteries are imaged.
A coronary artery moves fastest during systole and is best filled during diastole, so coronary
computed tomography angiography (CCTA) synchronizes acquisition to the electrocardiogram and
reconstructs the images in mid-diastole (Section~\ref{sec:ccta}).
The phase that best perfuses the arteries is also the phase in which they are most nearly still.

\begin{figure}[t]
  \centering
  \includegraphics[width=0.72\textwidth]{external/heart_wall.jpg}
  \caption{The coronary arteries run outside the muscle they supply. The muscular
    \emph{myocardium} lies between the endocardium and the epicardium, beneath which the arteries
    run; its inner and outer thirds are the subendocardium and subepicardium, which differ in both
    compression and oxygen extraction. Figure by OpenStax College, CC BY 3.0, via Wikimedia
    Commons.}
  \label{fig:heart-wall}
\end{figure}

\subsection{The coronary arterial tree}
\label{sec:tree}

Directions on the heart are given as \emph{anterior} (toward the front of the chest),
\emph{posterior} (toward the back), \emph{lateral} (toward the side) and \emph{inferior} (the
underside, which rests on the diaphragm).
Along a vessel, \emph{proximal} means near its origin at the aorta and \emph{distal} means far
from it, toward the small branches.

The coronary arteries arise from two of the three bulges in the aortic wall known as the aortic
sinuses \cite{devos2024coronary}.
The opening of each artery is its \emph{ostium}, and each ostium is the root of a tree: the left
coronary ostium opens from the left sinus and the right from the right sinus, which is why the two
systems are described separately here and stored as separate structures in ImageCAS-X
\cite{bransby2026imagecasx}.

\subsubsection{The left system}

The left coronary artery begins as the left main artery (LM), a short trunk that divides into the
left anterior descending artery (LAD) and the left circumflex artery (LCX)
\cite{devos2024coronary}.
In the 789 cohort cases where it is labeled, that trunk has a median length of 8.5 mm
(interquartile range 6.0--11.6 mm, full range 1.7--30.4 mm), so the first and largest branching
point of the left tree sits within a centimeter of the aorta.
The LAD descends in the anterior \emph{interventricular sulcus}, the groove between the two
ventricles.
It supplies the anterior wall of both, together with the anterior two-thirds of the
\emph{interventricular septum}, the muscular wall dividing them from each other
\cite{devos2024coronary,obrien2007anatomy}.
It does so through septal branches that penetrate that wall and a variable number of diagonal
branches that run out across the surface \cite{devos2024coronary}.
The LCX curves around the left side in the left \emph{atrioventricular sulcus}, the groove between
the left atrium and ventricle, and supplies the lateral wall through one or more obtuse marginal
branches \cite{devos2024coronary}.

Both sets of branches are numbered by the order in which they leave the parent vessel rather than
by any fixed identity: D1 and D2 are simply the first and second diagonal branches off the LAD, and
OM1 and OM2 the first and second obtuse marginal branches off the LCX \cite{devos2024coronary}.
Prevalence falls with position, and the cohort shows it: D1 is labeled in 98.5\% of cases and D2 in
59.4\%, OM1 in 86.8\% and OM2 in 37.6\%.

\subsubsection{The right system}

The right coronary artery (RCA) runs in the right atrioventricular sulcus, the groove between the
right atrium and the right ventricle \cite{devos2024coronary}.
It gives off branches to the right ventricle and to the \emph{sinoatrial node}, the tissue that sets
the heart's rhythm, then usually terminates on the underside of the heart as the posterior
descending artery (PDA), which supplies the inferior third of the septum, and one or more posterior
lateral arteries (PLA) \cite{devos2024coronary,obrien2007anatomy}.

Those proximal right-sided branches are not represented in the ImageCAS-X label set, which names
only the RCA trunk and the posterior vessels \cite{bransby2026imagecasx}.
This is an annotation boundary rather than an anatomical claim, and it constrains what a comparison
of branch counts between the left and right trees can mean.

\subsubsection{Territories}

The clinical weight of a lesion follows from the territory beyond it rather than from the lesion
itself.
An occlusion of the proximal LAD puts at risk the anterior wall and the septum, a large fraction of
the left ventricle, whereas an occlusion far along an obtuse marginal puts at risk a limited region of
the lateral wall \cite{devos2024coronary}.
Distance from the ostium is therefore a risk-relevant geometric quantity, a point returned to in
Section~\ref{sec:geometry}.

\begin{figure}[t]
  \centering
  \includegraphics[width=0.9\textwidth]{external/coronary_arteries.pdf}
  \caption{The named arteries of both systems, on the surface of the heart, seen from the front.
    Two naming differences from the convention used here: the circumflex artery is labeled Cx
    rather than LCX, and the left marginal branches M1 and M2 are the \emph{obtuse marginal}
    branches OM1 and OM2, named for the rounded left border along which they run. The right
    marginal artery (RMA) is drawn here but is not among the ImageCAS-X labels. Figure by Patrick
    J. Lynch, medical illustrator; derivative work: Fred the Oyster; Mikael H\"aggstr\"om, M.D,
    CC BY-SA 3.0, via Wikimedia Commons.}
  \label{fig:coronary-arteries}
\end{figure}

\begin{figure}[t]
  \centering
  \includegraphics[width=\textwidth]{anatomy_ct.png}
  \caption{The left main divides within millimeters of the aorta, and the annotation locates it to
    the voxel. Axial computed tomography (CT) slice of case 7 at the level of the coronary ostia,
    with the delivered
    segmentation overlaid in the per-artery colors used throughout. The ostium marker is the
    centerline's own \texttt{start\_points} flag, converted from the coordinate frame of the
    centerlines into voxel indices, where it samples 627 Hounsfield units (HU), the density scale of
    a CT image. This slice cuts through the aortic root rather than the ventricles,
    so only the aorta is named. Generated with \texttt{scripts/make\_anatomy\_figure.py}.}
  \label{fig:anatomy-ct}
\end{figure}

\subsection{Coronary dominance}
\label{sec:dominance}

The two systems differ in which of them reaches the \emph{crux}, the point on the underside of the
heart where the atrioventricular and interventricular sulci meet.
That variation is called \emph{dominance}.
It is defined by the origin of the posterior descending artery \cite{allwork1987applied}: a heart
is right-dominant when the PDA arises from the RCA, left-dominant when it arises from the LCX, and
codominant when the posterior territory is shared.
Dominance is the largest source of normal topological variation in the coronary tree, which makes
it the natural variable by which to stratify any description of coronary anatomy.

Approximately 85\% of people are right-dominant, 8\% left-dominant and 7\% codominant
\cite{wu2024dominance}.
These proportions vary systematically with ancestry rather than randomly: among patients undergoing
CCTA, right dominance has been reported in 96\% of East Asian, 87\% of South Asian and 87\% of
European individuals \cite{wu2024dominance}.
A cohort figure is therefore interpretable only alongside the population it was drawn from.

The ImageCAS-X cohort is annotated as 729 right-dominant, 41 left-dominant and 30 codominant cases,
or 91.1\%, 5.1\% and 3.8\% of the 800 usable cases \cite{bransby2026imagecasx}.
The underlying scans were acquired at a single hospital in Guangdong, China
\cite{zeng2023imagecas}, and the observed proportion of right dominance falls between the European
and East Asian figures above.

Dominance is recorded twice in ImageCAS-X, independently, and the two records agree
\cite{bransby2026imagecasx}: it appears as an explicit annotation, and it is implied by which
arteries are labeled, since the posterior vessels arise from whichever system is dominant.
Table~\ref{tab:posterior} gives the cross-tabulation.
The right posterior vessels appear in essentially every right-dominant case and in no left-dominant
case, the left posterior vessels invert the pattern exactly, and codominant cases draw on both.

\begin{table}[t]
  \centering
  \caption{The artery labels and the dominance annotation are two independent records of the same
    anatomical fact. Prevalence of the posterior vessels by annotated dominance over the 800 usable
    cases, from \texttt{scripts/cohort\_numbers.py}; the right and left posterior vessels are
    mutually exclusive in the single-dominance groups.}
  \label{tab:posterior}
  \begin{tabular}{lrrrr}
    \hline
    Vessel & All (\%) & Right (\%) & Left (\%) & Co (\%) \\
    \hline
    R-PDA & 93.0 & 99.7 &  0.0 & 56.7 \\
    R-PLA & 92.0 & 99.2 &  0.0 & 43.3 \\
    L-PDA &  5.8 &  0.0 & 90.2 & 30.0 \\
    L-PLA &  8.8 &  1.0 & 90.2 & 86.7 \\
    \hline
  \end{tabular}
\end{table}

\begin{figure}[t]
  \centering
  \includegraphics[width=\textwidth]{05_dominance_comparison.png}
  \caption{The posterior vessels follow the dominant system, so dominance changes which arteries
    exist rather than how large they are. Right-, left- and codominant coronary trees from the
    cohort; the right posterior descending artery is absent from the left-dominant tree entirely.}
  \label{fig:dominance}
\end{figure}
